\section{Hosting benchmarks in detail}\label{app:hosted}

This appendix gives the models, methods and results behind Table~\ref{tab:hosted:summary} of Section~\ref{sec:hosted}. Headline values
are not repeated here; the table is their home.

\subsection{Conventions}\label{app:hosted:conv}

\paragraph{Witnesses.} Witness~T and witness~I share no code path with each other or with the host. Where the host counts set bits with a
lookup table, witness~T uses bit-parallel arithmetic. Where witness~T reads a determinant from a pivoted LU factorisation, witness~I uses a
Householder QR factorisation. The gate passes only when both witnesses pass.

\paragraph{Controls.} Each benchmark carries four controls.
\begin{itemize}[nosep]
\item $[H,\Sigma]=0$ is checked on a small, complete Hilbert space.
\item An injected fault that leaves the existence sector must be refused by both witnesses.
\item Two orthogonal states of the sector, that is two different contents, must receive the same verdict.
\item In B4, a fault that changes only the content must pass both witnesses.
\end{itemize}

\paragraph{Uncertainties.} Exact identities (closed forms, commutators) are checked to floating-point accuracy. An extrapolated number
carries, as $1\sigma$, the largest change over a family of alternative fits (windows and functional forms) fixed in the code. In B3 a second
fit uses these spreads as weights; the uncertainty it propagates is combined in quadrature. A comparison passes at
$\vert\text{deviation}\vert\le2\sigma$.

\subsection{B1: the Heisenberg chain}\label{app:hosted:b1}

\paragraph{Model.} $H=\sum_{i=1}^{N}\mathbf S_i\cdot\mathbf S_{i+1}$ acts on spins $1/2$ on a ring, for even $N$ from 8 to 24. The
syndrome $\Sigma=(S^z_{\rm tot})^2$ selects the sector $S^z_{\rm tot}=0$, with 2\,704\,156 states at $N=24$. The ground-state energy per
site tends to $e_0=\tfrac14-\ln2$ \cite{bethe1931,hulthen1938}.

\paragraph{Method.} Lanczos iteration gives the two lowest levels. At every $N$ the ground energy is checked against the Bethe ansatz, the
exact solution of this chain. With $M=N/2$ rapidities $u_j$ and quantum numbers $I_j=-\tfrac{M-1}2,\dots,\tfrac{M-1}2$, the Bethe equations
and the energy are
\[
N\arctan2u_j=\pi I_j+\sum_k\arctan(u_j-u_k),\qquad E_0=\frac N4-\sum_j\frac2{1+4u_j^2}.
\]
A Newton solution agrees with Lanczos to $1.3\times10^{-15}$ per site. For a critical chain, conformal invariance gives
$E_0(N)/N=e_0-\pi vc/(6N^2)+\dots$ \cite{blote1986}. Here $c=1$ and $v=\pi/2$, and the form is corrected by powers of $1/\ln N$
\cite{affleck1989}. We fit $a+bN^{-2}+dN^{-4}$ on $N\ge12$. The nine alternatives are the windows $N\ge8,10,14,16$, and all five windows
with an added $N^{-6}$ term.

\paragraph{Results.} All nine alternatives lie above the exact value. The logarithmic corrections predict this, since no power series in
$1/N^2$ absorbs them. The triplet gap $\Delta_N=E_1-E_0$ closes:
\begin{itemize}[nosep]
\item $N\Delta_N$ grows from 4.18 at $N=8$ to 4.39 at $N=24$;
\item the local exponent $-\dd\ln\Delta_N/\dd\ln N$ rises from 0.946 to 0.968;
\item a fit $\Delta_\infty+\beta_1/N+\beta_2/N^2$ on $N\ge12$ gives $\Delta_\infty=0.0024$, against $\Delta_{24}=0.183$.
\end{itemize}
The chain is gapless, and $N\Delta_N$ approaches $\pi^2/2$ only logarithmically.

\paragraph{Witnesses.} T counts the up spins of every register index that carries amplitude. I evaluates $\langle(S^z_{\rm tot})^2\rangle$
site by site. $[H,S^z_{\rm tot}]=0$ holds exactly on the full $2^{10}$-dimensional space. The $N=24$ ground state and the orthogonal
triplet pass. A single spin flip gives $\langle(S^z_{\rm tot})^2\rangle=1$ and is refused by both.

\subsection{B2: the transverse-field Ising ring at \texorpdfstring{$L=4000$}{L = 4000}}\label{app:hosted:b2}

\paragraph{Model.} $H=-\sum_iZ_iZ_{i+1}-g\sum_iX_i$ acts on a ring of $L=4000$ spins, at $g=1$ (critical), $0.5$ and $1.5$. The
Jordan--Wigner transformation maps the spins to fermions. With the Majorana operators $a_j=(\prod_{l<j}X_l)Z_j$ and
$b_j=(\prod_{l<j}X_l)Y_j$, one has $X_j=ia_jb_j$ and $Z_jZ_{j+1}=ib_ja_{j+1}$. The bond that closes the ring carries the parity
$P=\prod_iX_i$. For $P=+1$ it becomes an antiperiodic fermion bond \cite{lieb1961,pfeuty1970}, and that sector is the existence sector.
There, and only there, $H=i\sum_{jk}A_{jk}a_jb_k$, with $A_{jj}=-g$, $A_{j+1,j}=1$, $A_{1L}=-1$ and all other entries zero. At $L=10$ this
reproduces the lowest even-parity eigenvalue of the full spin Hamiltonian to $1.8\times10^{-14}$.

\paragraph{Method.} The singular-value decomposition $A=U\,{\rm diag}(s_m)V^{\mathsf T}$ gives the ground energy $E_0=-\sum_ms_m$. The
ground state is Gaussian, with correlation matrix $G_{jk}=\langle ia_jb_k\rangle=-(UV^{\mathsf T})_{jk}$. A block of $l$ consecutive spins
has entropy
\[
S(l)=\sum_mh\bigl(\tfrac{1+\nu_m}2\bigr),\qquad h(p)=-p\ln p-(1-p)\ln(1-p),
\]
where $\nu_m$ are the singular values of $G$ restricted to the block \cite{peschel2003,vidal2003}. Each coupling costs one SVD of a
$4000\times4000$ matrix, about 27\,s on one core.

\paragraph{Results.} $E(m)$ and $K(k)$ denote the complete elliptic integrals of the second kind (parameter $m$) and of the first kind
(modulus $k$).
\begin{itemize}[nosep]
\item[(a)] The ground energy equals the finite-$L$ closed form $-\sum_k\sqrt{1+g^2-2g\cos k}$, with $k=2\pi(n+\tfrac12)/L$, at all three
couplings. At $g=0.5$ and $1.5$, $E_0/L$ equals the thermodynamic limit $e_\infty=-\tfrac2\pi(1+g)E(4g/(1+g)^2)$ to $4\times10^{-16}$
\cite{pfeuty1970}.
\item[(b)] At $g=1$, $E_0/L-e_\infty=-3.27\times10^{-8}$. Read as the Casimir term $-\pi vc/(6L^2)$ with $v=2$, it gives the first value of
$c$ in Table~\ref{tab:hosted:summary}. Its uncertainty is floating-point error, and the exact expansion predicts
$\tfrac12+7\pi^2/(480L^2)=0.5000000090$.
\item[(c)] At $g=1$ we fit the Calabrese--Cardy form for a ring \cite{calabrese2004}, $S(l)=\tfrac c3\ln[\tfrac L\pi\sin\tfrac{\pi l}L]+c_1'$,
on $100\le l\le2000$. This gives the second value of $c$, with $c_1'=0.478558$. The uncertainty is the spread over windows that start at
50, 200 or 400 or end at 1000.
\item[(d)] At $g=0.5$ and $1.5$ the entropy saturates, $S(L/2)-S(L/8)=2\times10^{-12}$: an area law. The saturated values match the
corner-transfer-matrix prediction. The corner transfer matrix is Baxter's transfer matrix for a quadrant of the two-dimensional lattice;
its spectrum gives the entanglement levels of a half-infinite chain. These are $(2j+1)\epsilon$ for $g>1$ and $2j\epsilon$ for $g<1$,
with $\epsilon=\pi K(k')/K(k)$ and $k=\min(g,1/g)$ \cite[Eqs.~4.30--4.34]{calabrese2004}. A block has two cuts, so $S=2S_{\rm half}$ for
$g>1$. For $g<1$ the even-parity ground state is a cat state whose $\ln2$ the two cuts share, so $S=2S_{\rm half}-\ln2$. We derived this
before the comparison.
\end{itemize}

\paragraph{Witnesses.} The syndrome is the parity, $\langle P\rangle=\det G$, together with purity ($G$ orthogonal). T reads $\det G$ from
a pivoted LU factorisation: parity $+1$, with $\vert\ln\vert\det G\vert\vert<10^{-13}$. I reads it from a Householder QR, using the
reflector count and the signs of the diagonal of $R$, and finds $\max\vert GG^{\mathsf T}-\mathbb1\vert\le7\times10^{-15}$. A spin flip
$Z_1$ reverses the first row of $G$; both witnesses read $-1$ and refuse.

\subsection{B3: the Schwinger model}\label{app:hosted:b3}

\paragraph{Model.} Quantum electrodynamics in 1+1 dimensions, with one fermion of mass $m$ and charge $g$, confines. At $m=0$ it is a free
boson of mass $g/\sqrt\pi$, the vector particle \cite{schwinger1962}. We use Kogut--Susskind staggered fermions \cite{kogut1975}: one
fermion component per site, the two components of the continuum spinor sitting on alternate sites. The lattice has $N$ sites, open
boundaries and zero background field. Gauss's law reads $E_n-E_{n-1}=Q_n$, with $Q_n=\tfrac12(\sigma^z_n+(-1)^n)$ the charge on site $n$ and
$E_n$ the integer electric flux on the link from site $n$ to $n+1$. It fixes $E_n=\sum_{k\le n}Q_k$ and leaves the spin Hamiltonian
\cite[Eq.~2.6]{banuls2013}
\begin{equation}\label{eq:hosted:schwinger}
W=\frac{2H}{ag^2}=x\sum_{n=0}^{N-2}\bigl(\sigma^+_n\sigma^-_{n+1}+{\rm h.c.}\bigr)+\sqrt x\,\frac{m_{\rm lat}}g\sum_{n=0}^{N-1}\bigl(1+(-1)^n
\sigma^z_n\bigr)+\sum_{n=0}^{N-2}E_n^2,
\end{equation}
with lattice spacing $a$ and $x=(ag)^{-2}$. The vector mass is $M/g=(W_1-W_0)/(2\sqrt x)$, from the two lowest levels, taken to
$N\to\infty$ and then to $x\to\infty$. We run two lattice masses: the naive $m_{\rm lat}=0$, and $m_{\rm lat}=-g^2a/8$. The discrete chiral
symmetry of the lattice model identifies the second as the massless point \cite{dempsey2022}. The shift vanishes as $a\to0$, so both
lattices have the same continuum limit.

\paragraph{Existence and witnesses.} The existence sector is fixed by Gauss's law at every site, with $E_{-1}=0$ (no background) and
$E_{N-1}=0$ (no field beyond the last site). Together these imply total charge zero. The host keeps an integer link register
$E_0,\dots,E_{N-2}$ for every basis configuration.
\begin{itemize}[nosep]
\item T checks $\mathcal G_n\equiv E_n-E_{n-1}-Q_n=0$ at every site of every configuration carrying amplitude. I evaluates
$\langle\sum_n\mathcal G_n^2+Q_{\rm tot}^2\rangle$.
\item The ground and vector states at $N=22$ pass. A charge without its field (one spin flipped) and a flux line without charges (one
link raised) give the value 2 and are refused by both.
\item The elimination itself is checked with explicit integer links at $N=6$ (1\,075\,648 states; $\vert E_n\vert\le3$ is exact there).
$[H,\mathcal G_n]=0$ for all $n$, and the 20 states with all $\mathcal G_n=0$ reproduce the spectrum of \eqref{eq:hosted:schwinger} to
$2.5\times10^{-14}$.
\end{itemize}

\paragraph{Method.} Lanczos iteration runs for $N=8,\dots,22$ (705\,432 states at $N=22$) at $x=0.5,1,1.5,2,3,4$.
\begin{itemize}[nosep]
\item \emph{$N\to\infty$ at each $x$.} We fit $M_N/g=w+\alpha N^{-2}+\beta N^{-3}$ on $N\ge12$. The $N^{-2}$ term is the kinetic energy
$\pi^2/(2M\ell^2)$ of the vector particle in a box of length $\ell$; the $N^{-3}$ term absorbs a boundary shift of $\ell$. The alternatives
are $N\ge14$, $N\ge16$, $\alpha N^{-2}+\beta N^{-4}$, and $\alpha N^{-2}$ alone on the last four sizes.
\item \emph{Which $x$ enter.} Only $x$ with $N_{\max}/\sqrt x\ge10$, a box at least ten units of $1/g$ long. All six qualify.
\item \emph{$x\to\infty$.} We fit $w(x)$ with a quadratic in $y=1/\sqrt x=ag$, weighted by the inverse squared $N$-extrapolation
errors. The alternatives are linear on $x\ge1.5$, quadratic on $x\ge1$, cubic, and, for the shifted lattice only, quadratic without the
linear term.
\end{itemize}

\begin{table}[htb]\centering\footnotesize
\caption{Schwinger model at $m=0$: vector mass $M/g$ at each $x=1/(ag)^2$, extrapolated to $N\to\infty$ (fit-family spread in units of
the last digit), and at $N=22$ \tagS~\chk{hosted:B3}. Exact value: $1/\sqrt\pi=0.5641896$.}\label{tab:hosted:schwinger}
\begin{tabular}{@{}llllll@{}}
\toprule
& & \multicolumn{2}{c}{shifted, $m_{\rm lat}=-g^2a/8$} & \multicolumn{2}{c}{naive, $m_{\rm lat}=0$}\\
\cmidrule(lr){3-4}\cmidrule(lr){5-6}
$x$ & $ag$ & $N\to\infty$ & $N=22$ & $N\to\infty$ & $N=22$\\
\midrule
0.5 & 1.41 & 0.5835(6) & 0.5892 & 0.8955(4) & 0.8986\\
1 & 1.00 & 0.5648(15) & 0.5784 & 0.7848(10) & 0.7939\\
1.5 & 0.82 & 0.5620(24) & 0.5836 & 0.7422(18) & 0.7580\\
2 & 0.71 & 0.5616(35) & 0.5910 & 0.7180(28) & 0.7405\\
3 & 0.58 & 0.5620(60) & 0.6067 & 0.6900(49) & 0.7259\\
4 & 0.50 & 0.5627(87) & 0.6222 & 0.6740(69) & 0.7228\\
\bottomrule
\end{tabular}
\end{table}

\paragraph{Results.} Table~\ref{tab:hosted:schwinger} gives the values at each $x$.
\begin{itemize}[nosep]
\item On the shifted lattice the continuum uncertainty combines a propagated $N$-extrapolation uncertainty of 0.016 with a continuum
fit-family spread of 0.035. On the naive lattice the two parts are 0.012 and 0.028.
\item The per-$x$ errors are fit-family spreads, strongly correlated between $x$. The main continuum fits have $\chi^2=0.04$ (shifted) and
0.18 (naive) for 3 degrees of freedom. With $\chi^2/{\rm dof}\ll1$, the deviations in units of $\sigma$ are consistency checks, not
significances.
\item The continuum alternatives on the shifted lattice give 0.562 (linear, $x\ge1.5$), 0.578 (quadratic, $x\ge1$), 0.570 (cubic) and
0.550 (no linear term).
\item The finite-size form passes an out-of-sample test. At $x=4$ the fit on $N\le22$ predicts the $N=24$ gap as 0.613467; the computed
value is 0.613388.
\end{itemize}

\paragraph{Precision limits.}
\begin{itemize}[nosep]
\item Matrix product states (MPS) give $0.56421(9)$ \cite{banuls2013}. The density-matrix renormalisation group (DMRG) gives $0.5642(2)$
\cite{byrnes2002}, as tabulated in \cite{banuls2013}. Their uncertainties are 200 to 400 times smaller than ours.
\item The MPS study uses $x$ from 5 to 600 with $N\ge20\sqrt x$ \cite{banuls2013}. Exact diagonalisation stops at $N=22$, hence at
$x\le4$.
\item The weighted continuum fit is dominated by the coarsest lattice, $x=0.5$. Its $N$-extrapolation is the most precise, but a
quadratic in $ag$ describes its lattice artefacts worst. This pulls the central value up and sets the 7\% error.
\item The shifted values at $1.5\le x\le4$ lie within 0.5\% of $1/\sqrt\pi$ \cite{dempsey2022}. They are not continuum values, and we
claim nothing from them.
\item We did not attempt the massive model, whose binding energies are known to $5\times10^{-5}$ \cite{banuls2013}, or the chiral
condensate. A continuum error of several per cent cannot test either.
\end{itemize}

\subsection{B4: Shor's period finding}\label{app:hosted:b4}

\paragraph{Circuit.} To factor $N$ we use a counting register of $t=2n$ qubits and a work register of $n=\lceil\log_2N\rceil$ qubits,
held as an exact state vector \cite{shor1997}. The run has three steps.
\begin{enumerate}[nosep]
\item The counting register starts in uniform superposition, and the work register in $\ket1$.
\item Controlled $U^{2^j}$ follow, with $U\ket y=\ket{ay\bmod N}$ for $y<N$ and $U\ket y=\ket y$ otherwise.
\item The inverse quantum Fourier transform is applied as the exact discrete Fourier transform. It agrees with the explicit gate circuit
to $1.8\times10^{-16}$.
\end{enumerate}
$U$ permutes $\{y<N\}$, so the existence sector (unit norm, work register $y<N$) is invariant. An outcome $k$ gives the period $r$ as
the first continued-fraction convergent denominator $q<N$ of $k/2^t$ with $a^q\equiv1$. It gives a factor $\gcd(a^{r/2}-1,N)$ when $r$
is even and $a^{r/2}\not\equiv-1\pmod N$; this is one run, with no repetition. We host $N=15,21,33,35,39,51,55,57$ with every base $a$
coprime to $N$, and $N=143$ (24 qubits) with $a=2,3,5,7$: 192 states in all.

\paragraph{Results.} The simulated distribution is compared with Shor's analytic distribution \cite{shor1997},
\[
P(k)=Q^{-2}\sum_{x_0<r}\Bigl\vert\sum_{j<m(x_0)}e^{-2\pi ijrk/Q}\Bigr\vert^2,\qquad Q=2^t,\quad m(x_0)=\lfloor(Q-1-x_0)/r\rfloor+1.
\]
Table~\ref{tab:hosted:shor} lists the register sizes and the success probabilities. The textbook failure appears where it should: for
$N=33$ and 57 with $a=2$ the period is found but $2^{r/2}\equiv-1$, so no factor follows.

\begin{table}[htb]\centering\footnotesize
\caption{Shor's period finding on an exact state vector \tagS~\chk{hosted:B4}. The simulated register has $3n$ qubits; Beauregard's
circuit needs $2n+3$ \cite{beauregard2003}. Exact single-run probabilities for $a=2$, and the mean over all hosted bases. $P$(factor) can
exceed $P$(period) because a recovered multiple of $r$ can also yield a factor.}\label{tab:hosted:shor}
\begin{tabular}{@{}rlrrrrrrr@{}}
\toprule
$N$ & factors & $3n$ & $2n+3$ & bases & $r$ ($a=2$) & $P$(period) & $P$(factor) & mean $P$(factor)\\
\midrule
15 & $3\times5$ & 12 & 11 & 7 & 4 & 0.500 & 0.500 & 0.429\\
21 & $3\times7$ & 15 & 13 & 11 & 6 & 0.331 & 0.331 & 0.211\\
33 & $3\times11$ & 18 & 15 & 19 & 10 & 0.398 & 0.000 & 0.220\\
35 & $5\times7$ & 18 & 15 & 23 & 12 & 0.331 & 0.331 & 0.303\\
39 & $3\times13$ & 18 & 15 & 23 & 12 & 0.331 & 0.331 & 0.303\\
51 & $3\times17$ & 18 & 15 & 31 & 8 & 0.500 & 0.500 & 0.484\\
55 & $5\times11$ & 18 & 15 & 39 & 20 & 0.393 & 0.393 & 0.320\\
57 & $3\times19$ & 18 & 15 & 35 & 18 & 0.329 & 0.000 & 0.180\\
143 & $11\times13$ & 24 & 19 & 4 & 60 & 0.265 & 0.265 & 0.232\\
\bottomrule
\end{tabular}
\end{table}

\paragraph{Witnesses.} T requires every amplitude with $y\ge N$ to be exactly zero and the norm to be one. I evaluates
$\langle\Pi_{y\ge N}\rangle=0$, with the projector applied as an operator. All 192 states pass.
\begin{itemize}[nosep]
\item A faulty multiplier that exchanges $\ket1$ and $\ket{2^n-1}$ moves amplitude out of the sector. Both witnesses refuse it in every
case.
\item A phase $Z$ on one work qubit changes the content but not the existence, and both witnesses pass it.
\end{itemize}

\subsection{The cost map}\label{app:hosted:cost}

One random layer consists of $n$ Haar-random (uniformly distributed) single-qubit unitaries and $n-1$ controlled-$Z$ gates. They act on
a dense state vector of $2^n$ amplitudes, 16 bytes each, on one core. The measured peak memory is the state plus a fixed 16\,MiB gate
buffer (Table~\ref{tab:hosted:cost}). The layer time divided by $n$ grows as $2^{1.068n}$ for $n=18$ to 26. The excess over $2^n$ comes from
cache and memory bandwidth. The rate is $6.9\times10^{-9}$\,s per qubit and amplitude for one layer, with the $n-1$ controlled-$Z$ gates
folded in. The extrapolations of Section~\ref{sec:hosted} use this rate.

\begin{table}[htb]\centering\footnotesize
\caption{Cost of a dense state-vector simulation, one random layer on one core \tagS~\chk{hosted:B4}.}\label{tab:hosted:cost}
\begin{tabular}{@{}rrrr@{}}
\toprule
$n$ & state & measured peak & seconds per layer\\
\midrule
10 & 16\,KiB & 54\,KiB & 0.0011\\
14 & 256\,KiB & 770\,KiB & 0.0024\\
18 & 4\,MiB & 10\,MiB & 0.032\\
22 & 64\,MiB & 80\,MiB & 0.63\\
24 & 256\,MiB & 272\,MiB & 3.7\\
26 & 1\,GiB & 1.02\,GiB & 15.5\\
\bottomrule
\end{tabular}
\end{table}

\subsection{Choices made after seeing the data}\label{app:hosted:choices}

\begin{enumerate}[nosep,label=(\roman*)]
\item In B1 the fit family and the $2\sigma$ rule were fixed after a prototype had shown the extrapolated value. Under a $1\sigma$ rule B1
would fail ($1.3\sigma$).
\item In B2 the entropy-fit windows were fixed after a prototype had shown $c=0.500000$ to six digits.
\item In B3 the $x$ grid and $N_{\max}=22$ were set from the cost of exploratory runs. The fit rules were written before the first
continuum fit (on $N\le18$) and were kept, with one exception. The alternative without a linear term was then dropped for the naive lattice,
whose mass shift is an $O(a)$ effect \cite{dempsey2022}. Keeping it widens the naive uncertainty from 0.031 to 0.075 (recomputed from
Table~\ref{tab:hosted:schwinger}); the comparison passes either way.
\item In B4 the faulty multiplier replaced an $X$ on the top work qubit as the fault. For $N=15$, $a=2$, that $X$ maps the orbit
$\{1,2,4,8\}$ to $\{9,10,12,0\}$, inside the existence sector.
\end{enumerate}

\subsection{Caveats}\label{app:hosted:caveats}

\begin{itemize}[nosep]
\item These are classical emulations of hosting. A classical host carries generic content only at a cost of order $2^n$.
\item The modular multiplications of B4 are exact permutation gates, not elementary-gate circuits. This changes gate counts but not
states. The registers use the textbook $t=2n$, not Beauregard's layout.
\item B1 stops at $N=24$, the memory limit of a sector code without translation symmetry.
\item Wall times are single-core on a shared machine and are reported, not checked. All other numbers are deterministic.
\end{itemize}
